\documentclass[../main.tex]{subfiles}

\begin{document}
\begin{CJK*}{UTF8}{gbsn}

\section*{Exercise 8}
Consider the full conformal prediction 
with $(X_1, Y_1), \cdots, (X_{n+1}, Y_{n+1})$ 
being exchangeable and a symmetric score function $s$,
prove that 

\begin{equation*}
    P \curly{Y_{n+1} \in \mathcal{C}(X_{n+1})} \leqslant 1 - \alpha + 
    \mathbb{E} \squared{\frac{1}{n+1}\sum_{i=1}^{n+1} \chi_{\{S_i = \hat{q}\}}}
\end{equation*}
where 

\begin{equation*}
    \hat{q} = \text{Quantile}\bracket{S_1, \ldots, S_{n+1}; 1- \alpha}.
\end{equation*}

Give an example showing that this bound is tighter than 

\begin{equation*}
    P \curly{Y_{n+1} \in \mathcal{C}(X_{n+1})} \leqslant \frac{\lceil (1-\alpha)(n+1) \rceil}{n+1}
    + \epsilon_{\text{tie}}
\end{equation*}
where $\epsilon_{\text{tie}}$
is the probability of $S_{n+1}$ 
being equal to some $S_j$ for $j \in \{1, \ldots, n\}$.

\paragraph{Solution}
As already discussed in the lecture and in the textbook,
$Y_{n+1} \in \mathcal{C}(X_{n+1})$ if and only if $S_{n+1} \leqslant \hat{q}$
and so 

\begin{equation*}
    P \curly{Y_{n+1} \in \mathcal{C}(X_{n+1})} =
    P \curly{S_{n+1} \leqslant \hat{q}} =
    P \curly{S_{n+1} < \hat{q}} + P \curly{S_{n+1} = \hat{q}}.
\end{equation*}
Given that $(S_1, \ldots, S_n,S_{n+1})$ is exchangeable,
$P \curly{S_{n+1} = \hat{q}} = P \curly{S_{i} = \hat{q}}$
for all $i \in \{1, \ldots, n\}$ and

\begin{equation*}
    P \curly{S_{n+1} = \hat{q}} 
    = \frac{1}{n+1} \sum_{i=1}^{n+1} P \curly{S_{i} = \hat{q}}
    = \mathbb{E} \squared{\frac{1}{n+1}\sum_{i=1}^{n+1} \chi_{\{S_i = \hat{q}\}}}.
\end{equation*}
The only thing left is to show that $P \curly{S_{n+1} < \hat{q}} \leqslant 1-\alpha$.
To do so, first observe that 
$S_{n+1} < \hat{q}$ if and only if $S_{n+1} < S_{(k)}$,
where $k = \lceil (1-\alpha)(n+1) \rceil$.
Thus:

\begin{equation*}
    P \curly{S_{n+1} < \hat{q}} 
    = P \curly{S_{n+1} < S_{(k)}} 
    \leqslant \frac{\lceil (1-\alpha)(n+1) \rceil - 1}{n+1}
    = \frac{\lfloor (1-\alpha)(n+1) \rfloor}{n+1} 
    \leqslant \frac{ (1-\alpha)(n+1) }{n+1} = 1-\alpha.
\end{equation*}
This completes the proof.

Finally, in case when the random variables $S_i$'s 
are all continuous random variables, 
$\epsilon_{\text{tie}} = P \curly{S_{n+1} = \hat{q}} = 0$ and 
evidently the bound that is worked out 
in this exercise is better.
In case when the random variables $S_i$'s 
are discrete and can only take very limited values,
it may happen that $S_1, \ldots, S_n$ has taken all 
possible values and $\epsilon_{\text{tie}} = 1$
so that the bound worked out here is really 
the only available bound.
////
\end{CJK*}
\end{document}